\clearpage
\section*{September 24, 2026: Estimating the size-and-splitting model from scratch}
\textit{Agent-drafted research entry.}

The September pilot reproduced the spike at 99 but not how parents split. It
predicted 4.4 percent of recent parents with two filings against 13.2
observed, 12.3 percent with three or more against 3.6, and almost no $99+99$
parents. Every alternative filing count drew its own cost of the same size, so a
parent that split spread across two to six filings. Jacob asked for a
splitting cost that rises with the number of buildings, so that pairs and
$99+99$ appear, and for a new estimation replacing
\path{fit_pure_notch_pilot}. It lives in the production task
\path{tasks/estimate_notch_model}.

\paragraph{The model.} It follows \path{framework_writeup.tex}. Each weighted
historical parent brings its preferred total $x$ and building count $J^0$.
Under the policy it chooses total units $m$ and buildings $J$ to minimize,
relative to its ordinary cost of 99 units, the size loss $((m-x)/99)^2$, a
burden on each building with 100 or more units (a jump $\kappa$ plus a kink
$\tau$ that raises the cost of each unit above 100), and a cost $c$ for each
building added beyond $J^0$. Parents differ in $c$, which is exponential with
mean $\sigma$. The historical organization is free, so the ordinary
multi-building parents of 2019--2022 stay as they were unless the threshold
gives them a reason to change. Layouts are free, sizes are exact integers, and
choice probabilities are computed exactly. The parent's own land does not
enter: consistent with the September 24 weights decision, zoning and borough
are fixed traits of the opportunity, and the size loss includes any change in
land. The weights match residential FAR and borough. The objective is the
sum of squared differences over the pilot's 45 cells (one, two or three-plus
buildings by 15 size bins), searched on a grid.

Before fitting, the task checks that without the burden every parent keeps
$x$ and $J^0$; that a single building shrinks to 99 exactly when
$x\le99(1+\sqrt\kappa)$ and is compressed by $1+\tau$ above; that minimum
burdens match exhaustive partitions for two and three buildings; that the
building menu never binds; that no parent consolidates; and that the exact
probabilities match simulated cost draws.

\paragraph{Estimates.} The jump is $\kappa=0.16$: crossing 100 costs a
single-building parent 16 percent of the ordinary cost of 99 units, so it
shrinks to 99 when its preferred size is below about 138. The mean splitting
cost is $\sigma=0.20$ per added building. The data choose no kink ($\tau=0$),
so the pure notch gives the same fit. Points within 10 percent of the best
objective span $\kappa$ from 0.12 to 0.30, $\tau$ from 0 to 0.02 and $\sigma$
from 0.16 to 0.50.

\begin{center}
\small
\begin{tabular}{lrrrr}
\hline
Share of recent parents (\%) & Observed & Benchmark & Model & Pilot \\
\hline
Exactly 99 units & 11.8 & 0.6 & 11.0 & 12.6 \\
Two buildings & 16.4 & 9.3 & 18.3 & 4.4 \\
Organized as $99+99$ & 3.0 & 0.0 & 1.4 & 0.2 \\
Three or more buildings & 4.3 & 3.3 & 8.6 & 12.3 \\
Above 300 units & 8.6 & 14.6 & 14.2 & 9.8 \\
Mean units & 139.5 & 170.9 & 169.0 & 149.9 \\
\hline
\end{tabular}
\end{center}

The benchmark is the reweighted 2019--2022 distribution. The pilot column is
the archived September 22 fit on the earlier panel, whose observed shares
differ (13.2 percent pairs, 2.9 percent $99+99$).

\begin{center}
\includegraphics[width=\textwidth]{../tasks/estimate_notch_model/output/pdf/model_fit.pdf}
\end{center}

The rising cost does what was intended. Pairs are now slightly
overpredicted, and the model produces half the observed $99+99$ parents,
against almost none before. Three-plus parents are still twice the observed
share, mostly above 300 units: when splitting is nearly free, a large
single-building parent keeps its size in as many buildings as it takes. The
exponential puts much of its mass near zero cost, so a distribution with less
mass there is the natural next specification.

\paragraph{Units.} Scaled to the 304 recent parents, the model implies 596
fewer proposed units. Holding each parent's building count fixed would lose
779, so splitting preserves 183. The raw difference between the weighted
historical total and the recent total is 9,555 units. Parents above 300 units
account for 10,743, more than all of it, with smaller sizes offsetting; the
model changes little there. The data reject compressing those
projects with a kink; the recent period simply has fewer large parents. With
lot-area weights the raw difference falls to 2,084 units and the model
explains 659 of them. Whether the missing large projects are a response to the
policy, such as assembling less land or not filing, or a change in available
sites is not something this model can settle.

\begin{center}
\small
\begin{tabular}{lrrrrr}
\hline
Specification & $\kappa$ & $\tau$ & $\sigma$ & Objective & Units lost \\
\hline
Main (jump and kink) & 0.16 & 0 & 0.20 & 0.0060 & 596 \\
Splitting cost $c, 3c, 6c$ & 0.16 & 0.02 & 0.32 & 0.0057 & 1,205 \\
Lot-area weights & 0.18 & 0 & 0.25 & 0.0024 & 659 \\
Common 180-day horizon & 0.16 & 0 & 0.16 & 0.0050 & 342 \\
Curvature 0.5 & 0.06 & 0 & 0.08 & 0.0059 & 689 \\
Curvature 2 & 0.60 & 0 & 0.79 & 0.0062 & 581 \\
Joint assessment & 0.03 & 0.20 & -- & 0.0099 & 6,950 \\
\hline
\end{tabular}
\end{center}

The jump and splitting cost are stable across weights and the common
follow-up horizon (187 recent parents). Curvature changes their scale, as it
should, not the fit. A splitting cost that rises faster per building fits
slightly better and doubles the units lost, because more parents shrink to
99 or $99+99$ (12.2 and 1.9 percent) instead of adding buildings; the unit
estimate depends on this shape.
Joint assessment of the whole parent fits much worse and produces no $99+99$
parents; splitting cannot help, so $\sigma$ is not identified.

\paragraph{Can the data tell the parameters apart?} Redrawing 200 recent
samples from the model at the estimate and re-estimating each, the jump falls
between 0.12 and 0.22 in 80 percent of draws, the kink is estimated at zero in
89 percent, and $\sigma$ lies between 0.13 and 0.32. If the true kink were 0.2,
its estimates would range from 0.05 to 0.6, with the jump anywhere from 0.02 to
1: a kink is weakly identified. These draws hold the historical sample fixed,
so they understate uncertainty.

\paragraph{Not yet done.} A splitting-cost distribution with less mass near
zero; inference that also resamples historical parents and weights; the
150-unit regime in Zones A and B; and the legal-assessment question for
observed splits.

\textit{Reproduction note.} Root \texttt{make estimate} runs
\path{tasks/estimate_notch_model}; results are in \path{estimates.csv},
\path{fit_moments.csv} and \path{recovery_summary.csv}.
